\section{Experiments}

\subsection{Experimental Scope}

Since SR-Workbench is currently positioned as a research substrate rather than a new symbolic regression algorithm, our experiments are designed to validate the platform along three axes: method coverage, dataset-schema validity, and end-to-end execution on a manageable pilot slice. We therefore do not treat the current results as a final leaderboard claim. Instead, the purpose of the experiments is to demonstrate that heterogeneous methods can be executed under a shared operating layer, that canonicalized datasets can be validated under a shared schema, and that the resulting workflow already supports OOD-aware reporting.

At the method level, the current registry integrates $10$ symbolic regression tools under a common interface: \texttt{QLattice}, \texttt{drsr}, \texttt{dso}, \texttt{e2esr}, \texttt{gplearn}, \texttt{iMCTS}, \texttt{llmsr}, \texttt{pyoperon}, \texttt{pysr}, and \texttt{tpsr}. These methods are mapped to $8$ conda environments in the current configuration, reflecting the fact that symbolic regression repositories often require incompatible dependency stacks.

\subsection{Dataset-Schema and Formula Validation}

We first validate the canonical dataset substrate itself. The current examples-style schema requires four aligned split files (\texttt{train.csv}, \texttt{valid.csv}, \texttt{id\_test.csv}, and \texttt{ood\_test.csv}) together with \texttt{metadata.yaml}, and optionally \texttt{formula.py} when analytical ground truth is available. Our validation script checks split presence, header consistency, target-feature alignment, declared sample counts, and optional formula consistency.

For datasets with known analytical formulas, the validator imports the declared formula callable, aligns its arguments with the feature ordering recorded in metadata, and compares formula outputs against stored target values. As a sanity check, the current implementation successfully validates the ground-truth formulas of \texttt{BPG0} and \texttt{BPG1} from the examples directory. On \texttt{BPG0}, the formula check yields split-level NMSEs on the order of $10^{-13}$ on train/validation/ID data and $10^{-9}$ on OOD data. On \texttt{BPG1}, the corresponding NMSEs are on the order of $10^{-15}$. These results indicate that the schema does not merely store formulas as annotations; it can also verify that a claimed analytical expression is numerically consistent with the dataset artifact.

\subsection{Pilot End-to-End Slice}

To assess the executable side of the platform, we report a pilot slice over three representative datasets already archived in the repository: \texttt{oscillator1}, \texttt{oscillator2}, and \texttt{stressstrain}. We use three methods that currently expose stable end-to-end execution paths through SR-Workbench: \texttt{pysr}, \texttt{llmsr}, and \texttt{drsr}. The reported results are taken from repository-archived \texttt{result.json} files generated through the unified execution path.

We emphasize again that these runs should be interpreted as \emph{substrate-validation case studies} rather than a definitive benchmark comparison. In particular, the current pilot slice is sufficient for showing that the platform can execute multiple heterogeneous methods under canonicalized datasets and produce aligned ID/OOD summaries, but it is not yet intended as the final frozen benchmarking protocol for the full paper.

Table~\ref{tab:pilot-slice} summarizes the current pilot results. Several useful patterns already emerge. First, the platform successfully reports both ID and OOD behavior under a common result structure, which is a basic but often missing requirement in practical symbolic regression evaluation. Second, OOD behavior is clearly more discriminative than ID behavior on the oscillator tasks. Third, LLM-assisted and classical methods can now be compared within the same execution substrate, even when they are backed by different repositories and runtime environments.

\begin{table}[t]
\centering
\scriptsize
\begin{tabular}{llrrrr}
\toprule
Method & Dataset & ID RMSE & OOD RMSE & ID NMSE & OOD NMSE \\
\midrule
\texttt{drsr} & \texttt{oscillator1} & $3.54\times 10^{-6}$ & $3.94\times 10^{-3}$ & $2.75\times 10^{-9}$ & $2.03\times 10^{-3}$ \\
\texttt{drsr} & \texttt{oscillator2} & $8.34\times 10^{-6}$ & $1.13\times 10^{-3}$ & $4.24\times 10^{-10}$ & $1.33\times 10^{-6}$ \\
\texttt{drsr} & \texttt{stressstrain} & $4.81\times 10^{-2}$ & $4.01\times 10^{-2}$ & $5.93\times 10^{-3}$ & $5.49\times 10^{-3}$ \\
\texttt{llmsr} & \texttt{oscillator1} & $7.01\times 10^{-4}$ & $2.75\times 10^{-2}$ & $1.08\times 10^{-4}$ & $9.91\times 10^{-2}$ \\
\texttt{llmsr} & \texttt{oscillator2} & $1.19\times 10^{-3}$ & $3.40\times 10^{-2}$ & $8.62\times 10^{-6}$ & $1.19\times 10^{-3}$ \\
\texttt{llmsr} & \texttt{stressstrain} & $5.56\times 10^{-2}$ & $4.99\times 10^{-2}$ & $7.90\times 10^{-3}$ & $8.51\times 10^{-3}$ \\
\texttt{pysr} & \texttt{oscillator1} & $1.99\times 10^{-3}$ & $6.17\times 10^{-2}$ & $8.68\times 10^{-4}$ & $4.97\times 10^{-1}$ \\
\texttt{pysr} & \texttt{oscillator2} & $1.55\times 10^{-1}$ & $1.11$ & $1.46\times 10^{-1}$ & $1.28$ \\
\texttt{pysr} & \texttt{stressstrain} & $5.23\times 10^{-2}$ & $4.88\times 10^{-2}$ & $7.00\times 10^{-3}$ & $8.13\times 10^{-3}$ \\
\bottomrule
\end{tabular}
\caption{Pilot end-to-end results from repository-archived runs executed through SR-Workbench. The current table is intended as a systems-validation slice rather than a final leaderboard.}
\label{tab:pilot-slice}
\end{table}

The pilot slice also highlights why an audited research substrate is valuable. On some tasks, especially \texttt{oscillator2}, the gap between ID and OOD behavior is large enough that purely ID-focused reporting would be misleading. On other tasks, such as \texttt{stressstrain}, multiple methods appear reasonably competitive on both ID and OOD splits, suggesting that later comparisons should include complexity and stability rather than relying on predictive error alone. These observations reinforce the design choice behind SR-Workbench: benchmark execution, schema design, and auditing logic should be part of the research contribution, not just hidden engineering details.

\subsection{Limitations of the Current Evaluation}

The present experiments are intentionally narrow. First, the pilot slice covers only a subset of integrated methods and tasks. Second, the current draft does not yet enforce a globally frozen cross-method compute budget for every supported backend. Third, expression-level normalization and complexity-aware ranking are not yet fully exploited in the reported tables. We therefore view the current experiments as a proof that the platform is already usable, while the next stage of the project will scale the evaluation protocol and strengthen the fairness guarantees required for broader benchmark claims.
